\clearpage
\section*{September 27, 2026: The scaled burden as the main model}
\textit{Agent-drafted research entry.}

Jacob made the scaled burden the main model for now: each parent's jump and
kink are multiplied by one lognormal scale with median 1, estimated by
likelihood with the unexplained share $\varepsilon$. It nests the single burden,
fits better than it or non-optimizers, and is where the reserved wage test
would enter. The least-squares single burden is kept as the robustness
specification \texttt{least\_squares} (formerly \texttt{main}); renaming
changed no result.

\paragraph{Bootstrap.} The 500 bootstrap samples are now drawn once, in
\path{draw_bootstrap_samples.R}, and shared by both bootstraps; the
single-burden bootstrap reproduces its earlier results exactly on them.
\path{bootstrap_scaled_burden.R} re-estimates the main model on its full grid
in every sample, and draw 0 reproduces the point estimate.

\begin{center}
\small
\begin{tabular}{lrr}
\hline
 & Estimate & Bootstrap 95\% interval \\
\hline
Median jump $\kappa$ & 0.18 & 0.06--0.50 \\
Median kink $\tau$ & 0 & 0--0.105 \\
Kink-to-jump ratio & 0 & 0--1.5 \\
Spread $s$ & 2 & 0--3 \\
Cost growth $\gamma$ & 2 & 1--3 \\
Mean splitting cost $\sigma$ & 0.20 & 0.10--0.50 \\
Unexplained share $\varepsilon$ & 0.015 & 0.001--0.05 \\
Units lost & 761 & 452--1,612 \\
\hline
\end{tabular}
\end{center}

The interval for units lost is narrower than the single burden's (396--2,197
by least squares, 337--2,981 by likelihood), and skewed up: the bootstrap
median is 872. Half the draws (53 percent) choose no kink. The growth of the
splitting cost is still not identified.

\paragraph{What the main model still misses.} Shares of the 278 recent parents
of at most 300 units; one percentage point is about three parents.

\begin{center}
\small
\begin{tabular}{lrrr}
\hline
 & Observed & Benchmark & Model \\
\hline
One building, 50--89 units & 49.3 & 46.8 & 45.6 \\
One building, 90--98 & 5.8 & 4.6 & 4.5 \\
One building, 99 & 12.9 & 0.4 & 10.9 \\
One building, 100--149 & 4.7 & 18.2 & 4.0 \\
One building, 150--300 & 7.9 & 20.7 & 10.1 \\
Two buildings, under 100 & 2.2 & 4.2 & 4.1 \\
Two buildings, 100--197 & 10.1 & 2.4 & 11.5 \\
$99+99$ & 3.2 & 0.0 & 3.4 \\
Two buildings, 199--300 & 0.7 & 0.8 & 0.5 \\
Three or more buildings & 3.2 & 1.9 & 5.4 \\
\hline
\end{tabular}
\end{center}

The largest misses are not responses to the threshold. Small parents with
several buildings, which the policy gives no reason to reorganize, are half as
common after the policy as before (2.2 against 4.2 percent in two buildings),
and the model, which leaves such parents as they were, overpredicts them and
the three-building parents. The difference persists with the common 180-day
horizon (2.4 against 5.1 percent of weighted parents under 100 units with
several buildings), so it is not late companion filings. Its mirror image is
the shortfall of small single buildings (55.1 against 50.1 percent under 99).
The splitting the policy causes, two buildings of 100--197 units, is close
(11.5 against 10.1). The model also puts too few parents at 99 and at 95--98
(3.2 percent observed against 1.0), since it sends every shrinking parent to
exactly 99, and keeps too many large single buildings (10.1 against 7.9).

\paragraph{Open questions.} In order of what they could change:
\begin{enumerate}
\item Whether the linkage is equally able to find small multi-building
parents in both periods. Historical parents also draw on evidence the recent
ones lack (Housing Database and parcel adjacency, reviewed links). If the
difference is linkage, the benchmark's building counts are not comparable.
\item Whether some affected parents are not built, or all project sizes
shifted between periods. Either would explain fewer large single buildings and
more small ones than the model predicts; exit would raise units lost, a common
shift would lower the policy's share of the change.
\item Whether parents that shrink land a few units below 99, which the model
could allow with a small landing error; this changes units lost little.
\end{enumerate}
